\scanpage{68}\mspage{C}
Reprenons alors $Z$ normal de dim $n\geq 2$, et
$X=Z-Y$. Plongeons $Z$ dans $P^r$, et le morphisme
induit $f:X\to P^r$. On prend $s=n-1$, de
sorte que $X_u$ \deleted{est une \uline{courbe}} \inserted{$=X_K\cap L_u=Z_K\cap L_u$ \deleted{\unsure{$=$}}, car $Y_K\cap L_u=\varnothing$,}
donc $X_u=Z_K\cap L_u$ est une \uline{courbe} \deleted{de
genre $g$} projective, d'ailleurs normale. Si $g$
est son genre, il résulte alors du corollaire
au lemme de Bertini que $\pi_1(X)$ admet
$2g$ générateurs topologiques, OK.

\uline{Conséquence}. Soit $K$ une extension
de type fini d'un corps $k$. Il
existe tjrs des modèles projectifs,
\deleted{\unsure{(donc|celui-ci} \unsure{normaux|normal}} à fortiori propres, et
normaux si on veut, \deleted{\unsure{donc|et} réguliers en
codim $\leq 1$, \inserted{Soit $X$ \unsure{de}} \unsure{la variété de|le modèle de} \unsure{des|du}
modèles \uline{réguliers} $U$ qui sont \deleted{\unsure{qui}}
\deleted{\unsure{par}} ``\uline{presque propres}'', i.e. complémentaires
d'un ens fermé de codim $\geq 2$ dans
un modèle propre $X$. D'ailleurs, \unsure{tous|ts|ttes} les
modèles $U$ comme dessus, \unsure{ordonnés|ordonnées} par}
% Small fork in the left margin of the cancelled paragraph.
\diagram{\begin{tikzcd}[ampersand replacement=\&,row sep=small,column sep=small]
{}\&X''\arrow[dl,dash]\arrow[dr,dash]\&{}\\
X\&{}\&X'
\end{tikzcd}}
% archive sheet number: 67
